% Part: first-order-logic
% Chapter: natural-deduction
% Section: proving-things

\documentclass[../../../include/open-logic-section]{subfiles}

\begin{document}

\iftag{FOL}
      {\olfileid{fol}{ntd}{pro}}
      {\olfileid{pl}{ntd}{pro}}

\olsection{اخذ (\usetoken{P}{derivation}) دیاں مثالاں}

\begin{ex}
آؤ !!{sentence} $(!A \land !B) \lif !A$ دا !!a{derivation} دیئے۔

پہلاں اسیں مطلوبہ نتیجہ !!{derivation} دے سب توں ہِٹھ لکھدے آں۔
\begin{prooftree}
\AxiomC{}
\UnaryInfC{$(!A\land !B) \lif !A$}
\end{prooftree}

اگلا سوال ایہ اے کہ ایس روپ دا !!a{sentence} کس قسم دے استنباط
توں نکل سکدا اے۔ نتیجے دا !!{main operator}~$\lif$ اے، ایس لئی
اسیں \Intro{\lif} دا قاعدہ لا کے نتیجے تک پہنچن دی کوشش کراں گے۔
نال نال ورتے ہوئے مفروضے لکھدے جاؤ تے استنباطی قاعدیاں دے ناں
لاؤ؛ ایس طرح ثبوت دے آخر وچ ویکھنا سوکھا رہندا اے کہ سارے مفروضے
!!{discharged} ہو چکے نیں یا نہیں۔
\begin{prooftree}
\AxiomC{$\Discharge{!A \land !B}{1}$}
\DeduceC{$!A$}
\DischargeRule{\Intro{\lif}}{1} 
\UnaryInfC{$(!A\land !B) \lif !A$}
\end{prooftree}

ہن مفروضے $!A \land !B$ توں $!A$ تک دے قدم بھرنے نیں۔ اتھے ساڈے
کول صرف اک منطقی رابط، $\land$، اے؛ ایس لئی $\land$ دا اخراجی
قاعدہ ورتنا پئے گا۔ ایہدے نال ہِٹھلا ثبوت ملدا اے:
\begin{prooftree}
\AxiomC{$\Discharge{!A \land !B}{1}$}
\RightLabel{\Elim{\land}}
\UnaryInfC{$!A$}
\DischargeRule{\Intro{\lif}}{1} 
\UnaryInfC{$(!A\land !B) \lif !A$}
\end{prooftree}
ہن ساڈے کول $(!A \land
!B) \lif !A$ دا درست !!{derivation} اے۔
\end{ex}

\begin{ex}
ہن آؤ $(\lnot !A \lor !B) \lif (!A \lif !B)$ دا !!a{derivation}
دیئے۔

پہلاں اسیں مطلوبہ نتیجہ !!{derivation} دے سب توں ہِٹھ لکھدے آں۔
\begin{prooftree}
\AxiomC{}
\UnaryInfC{$(\lnot !A \lor !B) \lif (!A \lif !B)$}
\end{prooftree}
ایہ نتیجہ دین والا منطقی قاعدہ لبھن لئی اسیں نتیجے دے منطقی
رابطے ویکھدے آں: $\lnot$، $\lor$ تے $\lif$۔ فی الحال ساڈا دھیان
$\lif$ دی پہلی ورتوں اُتے اے، کیوں جے اخذ دے آخری !!{sentence}
% PNBNTD-011: the ND tree ends in one sentence, not an end-sequent; retain the main-operator scope.
دا !!{main operator} ایہی اے۔ $\lnot$، $\lor$ تے $\lif$ دی دوجی
ورتوں کسے ہور رابطے دے دائرے وچ نیں، ایس لئی اوہناں نوں بعد وچ
سنبھالاں گے۔ سو اسیں \Intro{\lif} دے قاعدے توں شروع کردے آں۔
ایس دی درست ورتوں دا روپ ایہ ہووے گا:
\begin{prooftree}
\AxiomC{$\Discharge{\lnot !A \lor !B}{1}$}
\DeduceC{$!A \lif !B$}
\DischargeRule{\Intro{\lif}}{1}
\UnaryInfC{$(\lnot !A \lor !B) \lif (!A \lif !B)$}
\end{prooftree}
ہن اگے ودھن دے دو رستے نیں۔ یا اسیں ہِٹھوں اُپر ول کم جاری رکھ کے
\Intro{\lif} دی اک ہور ورتوں لبھیے، یا اُپروں ہِٹھاں ول کم کر کے
\Elim{\lor} دا قاعدہ لائیے۔ آؤ دوجا رستہ چنیے۔ مفروضے
$\lnot !A \lor !B$ نوں \Elim{\lor} دے سب توں کھبے مقدمے وجوں
ورتاں گے۔ \Elim{\lor} دی درست ورتوں وچ باقی دو مقدمے نتیجے
$!A \lif !B$ دے عین برابر ہونے چاہیدے نیں؛ البتہ ہر اک نوں
$\lnot !A \lor !B$ دے دو فصلی جزاں وچوں متعلقہ جز نوں مفروضہ
من کے وکھرا اخذ کیتا جا سکدا اے۔ ایس لئی ساڈا !!{derivation}
ایویں لگے گا:
\begin{prooftree}
\AxiomC{$\Discharge{\lnot !A \lor !B}{1}$}
\AxiomC{$\Discharge{\lnot !A}{2}$}
\DeduceC{$!A \lif !B$}
\AxiomC{$\Discharge{!B}{2}$}
\DeduceC{$!A \lif !B$}
\DischargeRule{\Elim{\lor}}{2}
\TrinaryInfC{$!A \lif !B$}
\DischargeRule{\Intro{\lif}}{1} 
\UnaryInfC{$(\lnot !A \lor !B) \lif (!A \lif !B)$}
\end{prooftree}

سجے پاسے دیاں دو ذیلی شاخاں وچ اسیں $!A \lif !B$ نوں !!{derive}
کرنا چاہندے آں؛ ایس لئی \Intro{\lif} ورتنا سب توں چنگا اے۔
\begin{prooftree}
\AxiomC{$\Discharge{\lnot !A \lor !B}{1}$}
\AxiomC{$\Discharge{\lnot !A}{2}, \Discharge{!A}{3}$}
\DeduceC{$!B$}
\DischargeRule{\Intro{\lif}}{3}
\UnaryInfC{$!A \lif !B$}
\AxiomC{$\Discharge{!B}{2}, \Discharge{!A}{4}$}
\DeduceC{$!B$}
\DischargeRule{\Intro{\lif}}{4}
\UnaryInfC{$!A \lif !B$}
\DischargeRule{\Elim{\lor}}{2}
\TrinaryInfC{$!A \lif !B$}
\DischargeRule{\Intro{\lif}}{1} 
\UnaryInfC{$(\lnot !A \lor !B) \lif (!A \lif !B)$}
\end{prooftree}

!!{derivation} دے باقی دو خالی حصیاں لئی سانوں $!B$ دے
!!{derivation}s چاہیدے نیں: وچکارلی شاخ وچ $\lnot !A$ تے $!A$
توں، تے سب توں سجی شاخ وچ $!A$ تے $!B$ توں۔
% PNBNTD-004: source calls the second subsidiary branch "on the left"; the tree makes it rightmost.
پہلاں وچکارلا حصہ لیندے آں۔ $\lnot
!A$ تے $!A$، \Elim{\lnot} دے دو مقدمے نیں:
\begin{prooftree}
\AxiomC{$\Discharge{\lnot !A}{2}$}
\AxiomC{$\Discharge{!A}{3}$}
\RightLabel{\Elim{\lnot}}
\BinaryInfC{$\lfalse$}
\DeduceC{$!B$}
\end{prooftree}
\FalseInt لا کے اسیں $!B$ نوں نتیجے وجوں حاصل کر سکدے آں
تے ایہ شاخ مکمل ہو جاندی اے۔
\begin{prooftree}
\AxiomC{$\Discharge{\lnot !A \lor !B}{1}$}
\AxiomC{$\Discharge{\lnot !A}{2}$}
\AxiomC{$\Discharge{!A}{3}$}
% PNBNTD-005: source intermediate label Intro(false) contradicts its Elim(not) rule and finished tree.
\RightLabel{\Elim{\lnot}}
\BinaryInfC{$\lfalse$}
\RightLabel{\FalseInt}
\UnaryInfC{$!B$}
\DischargeRule{\Intro{\lif}}{3}
\UnaryInfC{$!A \lif !B$}
\AxiomC{$\Discharge{!B}{2}, \Discharge{!A}{4}$}
\DeduceC{$!B$}
\DischargeRule{\Intro{\lif}}{4}
\UnaryInfC{$!A \lif !B$}
\DischargeRule{\Elim{\lor}}{2}
\TrinaryInfC{$!A \lif !B$}
\DischargeRule{\Intro{\lif}}{1} 
\UnaryInfC{$(\lnot !A \lor !B) \lif (!A \lif !B)$}
\end{prooftree}

ہن سب توں سجی شاخ ویکھو۔ اتھے ایہ سمجھنا ضروری اے کہ
!!{derivation} دی تعریف مفروضے \emph{خارج کرن دی اجازت دیندی اے}،
پر اوہناں نوں خارج کرنا \emph{لازمی نہیں بناؤندی}۔ دوجے لفظاں وچ،
جے اسیں $!B$ نوں مفروضیاں $!A$ تے $!B$ وچوں اک توں
اخذ کر لئیے، دوجے نوں ورتے بنا، تاں ایہ ٹھیک اے۔ $!B$ نوں خود $!B$ توں
!!{derive} کرنا بالکل سادہ اے: اکیلا $!B$ ایہو جیہا
!!a{derivation} اے تے کوئی استنباط نہیں چاہیدا۔ سو اسیں
مفروضہ~$!A$ ہٹا سکدے آں۔
\begin{prooftree}
\AxiomC{$\Discharge{\lnot !A \lor !B}{1}$}
\AxiomC{$\Discharge{\lnot !A}{2}$}
\AxiomC{$\Discharge{!A}{3}$}
\RightLabel{\Elim{\lnot}}
\BinaryInfC{$\lfalse$}
\RightLabel{\FalseInt}
\UnaryInfC{$!B$}
\DischargeRule{\Intro{\lif}}{3}
\UnaryInfC{$!A \lif !B$}
\AxiomC{$\Discharge{!B}{2}$}
\RightLabel{\Intro{\lif}}
\UnaryInfC{$!A \lif !B$}
\DischargeRule{\Elim{\lor}}{2}
\TrinaryInfC{$!A \lif !B$}
\DischargeRule{\Intro{\lif}}{1} 
\UnaryInfC{$(\lnot !A \lor !B) \lif (!A \lif !B)$}
\end{prooftree}
دھیان دیو کہ مکمل !!{derivation} وچ سب توں سجی \Intro{\lif}
والی استنباطی ورتوں حقیقت وچ کوئی مفروضہ خارج نہیں کردی۔
\end{ex}

\begin{ex}
ہن تک سانوں \FalseCl{} دے قاعدے دی لوڑ نہیں پئی۔ ایہ ایس گل
وچ خاص اے کہ ایہ اوہ مفروضہ وی خارج کرن دی اجازت دیندا اے جیہڑا
قاعدے دے نتیجے دا ذیلی-!!{formula} نہیں ہوندا۔ ایہ \FalseInt{}
دے قاعدے نال نیڑے جڑیا ہویا اے۔ اصل وچ \FalseInt{} دا قاعدہ،
\FalseCl{} دے قاعدے دی اک خاص صورت اے---اک منطق، ”وجدانیت پسند
منطق“ (intuitionistic logic)، وی اے جیہدے وچ صرف \FalseInt{}
دی اجازت اے۔ جدوں ہور کوئی طریقہ کم نہ کرے، \FalseCl{} آخری
چارہ اے۔ مثلاً فرض کرو اسیں $!A \lor \lnot !A$ نوں !!{derive}
کرنا چاہندے آں۔ ساڈی عام تدبیر $!A \lor \lnot !A$ نوں
$\Intro{\lor}$ لا کے !!{derive} کرن دی ہووے گی۔ پر ایس لئی
بنا کسے مفروضے دے یا $!A$ یا $\lnot !A$ نوں !!{derive} کرنا
پئے گا، تے ایہ نہیں ہو سکدا۔ ہن \FalseCl{} کم آؤندا اے:
\begin{prooftree}
  \AxiomC{$\Discharge{\lnot(!A \lor \lnot !A)}{1}$}
  \DeduceC{$\lfalse$}
  \DischargeRule{\FalseCl}{1}
  \UnaryInfC{$!A \lor \lnot !A$}
\end{prooftree}
ہن سانوں $\lfalse$ دا !!a{derivation}،
$\lnot(!A \lor \lnot !A)$ توں چاہیدا اے۔ کیوں جے $\lfalse$،
$\Elim{\lnot}$ دا نتیجہ اے، اسیں اوہ قاعدہ آزما سکدے آں:
\begin{prooftree}
  \AxiomC{$\Discharge{\lnot(!A \lor \lnot !A)}{1}$}
  \DeduceC{$\lnot !A$}
  \AxiomC{$\Discharge{\lnot(!A \lor \lnot !A)}{1}$}
  \DeduceC{$!A$}
  \RightLabel{\Elim{\lnot}}
  \BinaryInfC{$\lfalse$}
  \DischargeRule{\FalseCl}{1}
  \UnaryInfC{$!A \lor \lnot !A$}
\end{prooftree}
$\lnot !A$ دا !!a{derivation} لبھن دی ساڈی تدبیر
$\Intro{\lnot}$ لاون نوں آکھدی اے:
\begin{prooftree}
  \AxiomC{$\Discharge{\lnot(!A \lor \lnot !A)}{1}, \Discharge{!A}{2}$}
  \DeduceC{$\lfalse$}
  \DischargeRule{\Intro{\lnot}}{2}
  \UnaryInfC{$\lnot !A$}
  \AxiomC{$\Discharge{\lnot(!A \lor \lnot !A)}{1}$}
  \DeduceC{$!A$}
  \RightLabel{\Elim{\lnot}}
  \BinaryInfC{$\lfalse$}
  \DischargeRule{\FalseCl}{1}
  \UnaryInfC{$!A \lor \lnot !A$}
\end{prooftree}
اتھے اسیں $\lfalse$ سوکھ نال حاصل کر سکدے آں: $\Elim{\lnot}$
نوں مفروضے $\lnot(!A \lor \lnot !A)$ تے
$!A \lor \lnot !A$ اُتے لا کے۔ دوجا مقدمہ ساڈے نویں مفروضے $!A$
توں $\Intro{\lor}$ راہیں نکلدا اے:
\begin{prooftree}
  \AxiomC{$\Discharge{\lnot(!A \lor \lnot !A)}{1}$}
  \AxiomC{$\Discharge{!A}{2}$}
  \RightLabel{\Intro{\lor}}
  \UnaryInfC{$!A \lor \lnot !A$}
  \RightLabel{\Elim{\lnot}}
  \BinaryInfC{$\lfalse$}
  \DischargeRule{\Intro{\lnot}}{2}
  \UnaryInfC{$\lnot !A$}
  \AxiomC{$\Discharge{\lnot(!A \lor \lnot !A)}{1}$}
  \DeduceC{$!A$}
  \RightLabel{\Elim{\lnot}}
  \BinaryInfC{$\lfalse$}
  \DischargeRule{\FalseCl}{1}
  \UnaryInfC{$!A \lor \lnot !A$}
\end{prooftree}
سجے پاسے اسیں اوہی تدبیر ورتدے آں، پر اتھے $!A$ نوں
\FalseCl راہیں حاصل کردے آں:
\begin{prooftree}
  \AxiomC{$\Discharge{\lnot(!A \lor \lnot !A)}{1}$}
  \AxiomC{$\Discharge{!A}{2}$}
  \RightLabel{\Intro{\lor}}
  \UnaryInfC{$!A \lor \lnot !A$}
  \RightLabel{\Elim{\lnot}}
  \BinaryInfC{$\lfalse$}
  \DischargeRule{\Intro{\lnot}}{2}
  \UnaryInfC{$\lnot !A$}
  \AxiomC{$\Discharge{\lnot(!A \lor \lnot !A)}{1}$}
  \AxiomC{$\Discharge{\lnot !A}{3}$}
  \RightLabel{\Intro{\lor}}
  \UnaryInfC{$!A \lor \lnot !A$}
  \RightLabel{\Elim{\lnot}}
  \BinaryInfC{$\lfalse$}
  \DischargeRule{\FalseCl}{3}
  \UnaryInfC{$!A$}
  \RightLabel{\Elim{\lnot}}
  \BinaryInfC{$\lfalse$}
  \DischargeRule{\FalseCl}{1}
  \UnaryInfC{$!A \lor \lnot !A$}
\end{prooftree}
\end{ex}

\begin{prob}
ہِٹھ لکھیاں گلّاں وکھاؤن والے !!{derivation}s دیو:
\begin{enumerate}
\item $!A \land (!B \land !C) \Proves (!A \land !B) \land !C$.
\item $!A \lor (!B \lor !C) \Proves (!A \lor !B) \lor !C$.
\item $!A \lif (!B \lif !C) \Proves !B \lif (!A \lif !C)$.
\item $!A \Proves \lnot\lnot !A$.
\end{enumerate}
\end{prob}

\begin{prob}
ہِٹھ لکھیاں گلّاں وکھاؤن والے !!{derivation}s دیو:
\begin{enumerate}
\item $(!A \lor !B) \lif !C \Proves !A \lif !C$.
\item $(!A \lif !C) \land (!B \lif !C) \Proves (!A \lor !B) \lif !C$.
\item $\Proves \lnot(!A \land \lnot !A)$.
\item $!B \lif !A \Proves \lnot !A \lif \lnot !B$.
\item $\Proves (!A \lif \lnot !A) \lif \lnot !A$.
\item $\Proves \lnot(!A \lif !B) \lif \lnot !B$.
\item $!A \lif !C \Proves \lnot (!A \land \lnot !C)$.
\item $!A \land \lnot !C \Proves \lnot (!A \lif !C)$.
\item $!A \lor !B, \lnot !B \Proves !A$.
\item $\lnot !A \lor \lnot !B \Proves \lnot(!A \land !B)$.
\item $\Proves (\lnot !A \land \lnot !B) \lif\lnot(!A \lor !B)$.
\item $\Proves \lnot(!A \lor !B) \lif (\lnot !A \land \lnot !B)$.
\end{enumerate}
\end{prob}

\begin{prob}
ہِٹھ لکھیاں گلّاں وکھاؤن والے !!{derivation}s دیو:
\begin{enumerate}
\item $\lnot(!A \lif !B) \Proves !A$.
\item $\lnot(!A \land !B) \Proves \lnot !A \lor \lnot !B$.
\item $!A \lif !B \Proves \lnot !A \lor !B$.
\item $\Proves \lnot \lnot !A \lif !A$.
\item $!A \lif !B, \lnot !A \lif !B \Proves !B$.
\item $(!A \land !B) \lif !C \Proves (!A \lif !C) \lor (!B \lif !C)$.
\item $(!A \lif !B) \lif !A \Proves !A$.
\item $\Proves (!A \lif !B) \lor (!B \lif !C)$.
\end{enumerate}
(ایہناں ساریاں لئی $\FalseCl$~دا قاعدہ ضروری اے۔)
\end{prob}

\end{document}
